\begin{tcolorbox}[
  colback=yellow!7,
  colframe=blue!35!black,
  colbacktitle=yellow!75!orange,
  coltitle=blue!25!black,
  title=\MakeUppercase{Did You Know?},
  fonttitle=\large\bfseries\sffamily,
  toptitle=1.5mm,
  bottomtitle=1.5mm,
  boxrule=0.7pt,
  arc=0mm,
  left=4mm,
  right=4mm,
  top=2.5mm,
  bottom=2.5mm
]
\footnotesize
\noindent
\begin{minipage}[t]{0.48\linewidth}
  {\sffamily\bfseries\textcolor{blue!35!black}{THE 2.15 dB DIFFERENCE}}\par
  A half-wave dipole has approximately \(2.15\text{ dBi}\) gain. Therefore,
  \(0\text{ dBd}=2.15\text{ dBi}\).

  \vspace{0.18em}
  {\sffamily\bfseries\textcolor{blue!35!black}{GAIN REDIRECTS ENERGY}}\par
  Antenna gain does not create extra transmitter power; it concentrates the
  available radiation more strongly in selected directions.
\end{minipage}
\hfill
\begin{minipage}[t]{0.48\linewidth}
  {\sffamily\bfseries\textcolor{blue!35!black}{THE 3 dB RULE}}\par
  A \(3\text{ dB}\) increase represents approximately twice the power, while
  a \(3\text{ dB}\) decrease represents half the power.

  \vspace{0.18em}
  {\sffamily\bfseries\textcolor{blue!35!black}{FREQUENCY SETS THE SCALE}}\par
  Since \(\lambda=c/f\), antennas designed for higher frequencies are
  generally physically smaller than antennas designed for lower frequencies.
\end{minipage}
\end{tcolorbox}